\section{Repair Funnel and Failure Analysis (RQ4)}
\label{sec:engineering}
\subsection{The Four-Step Repair Funnel}
High review coverage does not guarantee a correct diagnosis or repair. To answer RQ4, Table~\ref{tab:engineering} follows four nested steps across all 387 admitted patches for each investigator and configuration.

\begin{table*}[t]
\caption{Four-step investigation and repair funnel across 387 admitted patches per configuration. SecPass and new fixes count only candidates within Step 4. Bold highlights supported assessments and C2's new fixes.}
\label{tab:engineering}
\centering\footnotesize
% Generated by paper/build_final_tables.py; FINAL_MANUSCRIPT_DATA.json
\begin{tabular}{@{}lrrrrrr@{}}
\toprule
\multicolumn{7}{@{}l}{\textbf{A. Fixed ten-case development cohort}} \\
Version & C0 S & C3 S & Reject. & Cite & Limit & Error \\
\midrule
v6 & 1 & 0 & 31/48 & 15/23 & 3/2 & 0/0 \\
v7 & 1 & 1 & 64/73 & 5/20 & 3/2 & 0/0 \\
v8 & 1 & 0 & 23/9 & 0/0 & 2/1 & 0/0 \\
v9 & 1 & 1 & 92/62 & 0/0 & 9/9 & 0/1 \\
v10 & 1 & 3 & 49/61 & 0/0 & 3/4 & 1/1 \\
v11 & 1 & 3 & 57/30 & 0/0 & 4/4 & 0/2 \\
v12 & 2 & 4 & 77/43 & 0/1 & 5/6 & 0/0 \\
\bottomrule
\end{tabular}
\par\smallskip
\begin{tabular}{@{}llrrrrrr@{}}
\toprule
\multicolumn{8}{@{}l}{\textbf{B. Full recorded funnel, 387 attempts per arm}} \\
Inv. & Arm & S1 & S2 & S3 & S4 & Sec & New \\
\midrule
D & C0 & 245 & 109 & 89 & 58 & 13 & 4 \\
D & C1 & 251 & 109 & 98 & 65 & 13 & 5 \\
D & C2 & 237 & 85 & 74 & 54 & 13 & 8 \\
D & C3 & 252 & 97 & 89 & 60 & 10 & 6 \\
\midrule
G & C0 & 213 & 15 & 14 & 11 & 3 & 2 \\
G & C1 & 225 & 32 & 28 & 20 & 8 & 6 \\
G & C2 & 223 & 29 & 27 & 18 & 8 & 7 \\
G & C3 & 235 & 23 & 21 & 17 & 5 & 4 \\
\bottomrule
\end{tabular}

\end{table*}

\begin{enumerate}
\item \textbf{Hypothesis on Target Symbol.} The investigator records a hypothesis on a symbol touched by the reference security patch. DeepSeek reaches this step in 237 to 252 cases per configuration, and Gemini in 213 to 235.
\item \textbf{Supported Assessment.} The investigator records an accepted supported assessment on a reference symbol. Counts fall to 85 to 109 for DeepSeek and 15 to 32 for Gemini. This records the investigator's assessment, not independent confirmation of the true flaw.
\item \textbf{Reference-File Edit.} A case with supported localization also contains an edit in a reference file. This does not require that the edit address the same symbol or hypothesis. Counts range from 74 to 98 for DeepSeek and 14 to 28 for Gemini.
\item \textbf{Nonempty Retained Patch.} The case ends with a nonempty candidate patch, leaving 54 to 65 cases for DeepSeek and 11 to 20 for Gemini. Retention alone is not a functional or security verdict. The final column reports official security outcomes for this retained subset.
\end{enumerate}

Each step includes only cases that reached the preceding steps, so the table locates where a candidate leaves the observed workflow. Its last column counts security passes among retained candidates and identifies which are new conversions. It therefore answers a different question from Table~\ref{tab:results}, which includes the final outcomes of all evaluated patches, including unchanged outputs.

\subsection{Drop-off 1. False Refutations During Diagnosis}
The largest loss occurs between forming a hypothesis and recording support. Decision ledgers show premature refutations, wrong-property hypotheses, inconclusive assessments, and turn exhaustion. Some investigators inspect partial checks or helper calls and conclude that the boundary is already protected. For example, a RestrictedPython investigation treats existing underscore restrictions as sufficient despite missing restrictions on inspection attributes. Other investigations examine an unrelated property after finding the correct function. An Aiohttp session hypothesis can concern dictionary copying while the actual contract concerns expiry. Because the funnel also includes initially secure outputs, an accepted refutation is not automatically a false all-clear.

\subsection{Drop-off 2. Turn Exhaustion Before Validation}
Late edits can exhaust the turn budget before public validation finishes. Validation slots and required probe reruns can also prevent retention. A separate development recovery check evaluates seventy saved candidates discarded for protocol completion issues. Of these, 68 hit turn or validation-slot limits and two lacked a probe rerun. Offline tests find 39 functional passes, eighteen failures, and thirteen unresolved executions. Nine pass security, comprising eight preserved successes and one new conversion. These recovery results remain separate from the reported campaign scores. They motivate a reserved final validation opportunity after editing, without establishing that every discarded patch was correct or lost to scheduling.

\subsection{Drop-off 3. Incomplete Fixes and Contract Mismatches}
\begin{table}[t]
\caption{Manual inspection of fifteen retained patches with reference-symbol edits, verified functional passes, and failed security tests. Categories describe this selected diagnostic sample.}
\label{tab:taxonomy}
\centering\footnotesize
% Qualitative inspection of 15 retained candidates, not a population estimate.
\begin{tabularx}{\columnwidth}{@{}>{\raggedright\arraybackslash}p{0.25\columnwidth}r>{\raggedright\arraybackslash}X@{}}
\toprule
Failure category & Count & Repair attempt and remaining defect \\
\midrule
\textbf{Right bug class, incomplete fix} & \textbf{9} & \textbf{A plausible fix leaves a bypass or incomplete boundary.} Streamlit misses symlink escape. JWCrypto checks size after decompression. Airflow Hive skips inherited schema. Waitress misses framing variants. Flask-Admin leaves a backslash URL bypass. \\
\addlinespace
Wrong security property & 4 & Aiohttp checks session keys rather than expiry. FastAPI limits body size rather than media-type handling. Other candidates leave route anchoring or username normalization unresolved. \\
\addlinespace
Right condition, wrong failure behavior & 2 & NumPy raises \code{IndexError} instead of \code{ValueError} on an empty axis. Markdown-it-py still raises on malformed-byte decoding. \\
\bottomrule
\end{tabularx}

\end{table}
Even a retained patch that passes functional tests can leave the security contract broken. We inspect fifteen such tasks outside the diagnostic pilot by comparing candidate diffs, reference patches, and official failing assertions. Fourteen use DeepSeek and one uses Gemini.

Table~\ref{tab:taxonomy} shows nine incomplete repairs of the relevant vulnerability class. Examples include lexical path normalization that misses symlinks in Streamlit, size checks after decompression in JWCrypto, and validation that omits inherited connection settings in Airflow Hive. Four patches address the wrong security property, such as session-key handling instead of expiry in Aiohttp, or body-size limits instead of media-type parsing in FastAPI. Two violate the required failure behavior. NumPy raises the wrong exception for an empty axis, while Markdown-it-py leaves malformed-byte handling incorrect. These failures require checking the full input and error contract, not merely the presence of a new guard.

\subsection{Operational and Provider Failures}
Gemini records 205 repair errors. These comprise 153 \code{blockReason: OTHER} responses, 45 JSON decoding errors, six HTTP 500 errors, and one other value error. DeepSeek records 107 errors, including 63 explicit stop interruptions, 32 HTTP 429 responses, nine deadlines, two value errors, and one connection error. The recorded labels distinguish provider responses, parsing failures, and run interruptions without assigning unverified causes. Final candidate evaluations remain separate from these investigation errors.

\rqfinding{RQ4 Finding}{Candidates are lost during diagnosis, validation scheduling, and semantic repair. The funnel distinguishes these losses from operational errors and from the final test outcome.}

\section{Industrial Case Study. Structural Biology Pipelines (RQ5)}
\label{sec:practice}
Silent-fix research motivates retaining repair purpose beyond code~\cite{colefunda}. We evaluate maintenance independently of SusVibes through ten feature tasks in five scientific repositories, separating repair memory from boundary feedback.

\begin{table*}[t]
\caption{Independent biology maintenance study: two tasks per repository. Repository rows report final FuncPass/SecPass/JointPass after at most one retry; aggregate first and final outcomes follow. FuncPass requires ordinary and feature checks; SecPass requires saved and extended boundary checks, including reliability contracts for DockQ and AutoDock-Vina. Totals have denominator ten.}
\label{tab:bio}
\centering\footnotesize
% Evidence: silent_fix_feature_v1/{cases,continuations}; repository counts:
% paper/NOSTRUM_BARCELONA_PROOFREAD_PACK.md, original pinned scan inventory.
\setlength{\tabcolsep}{3pt}
\renewcommand{\arraystretch}{1.12}
\begin{tabularx}{\textwidth}{@{}l r >{\raggedright\arraybackslash}p{0.17\textwidth} >{\raggedright\arraybackslash}X >{\raggedright\arraybackslash}X c c@{}}
\toprule
Repository & Alerts & Target symbol & Repaired component defect & Feature evolution & \multicolumn{2}{c}{Joint success}\\
\cmidrule(lr){6-7}
 & & & & & No Memory & Memory+Gate\\
\midrule
DockQ & 10 & \code{load_PDB}
& Handle leak on parser failure/fallback
& Extension-first CIF/PDB parsing, optional format return
& Yes & Yes\\
haddock3 & 247 & \code{extract_files_flat}
& Extraction escape through destination symlink
& ZIP and TAR support, optional nested paths
& No$^{a}$ & Yes\\
ProteinMPNN & 43 & \code{StructureDataset}\newline\code{.__init__}
& Unvalidated record names used in output paths
& Minimum-length and missing-coordinate filters
& Yes & Yes\\
RFdiffusion & 49 & \code{Sampler.}\newline\code{load_checkpoint}
& Unrestricted pickle loading under pinned runtime$^{b}$
& Legacy-checkpoint exception handling and CPU fallback
& No$^{c}$ & Yes, retry\\
AutoDock-Vina & 60 & \code{Vina.write_poses}
& Overwrite branch bypasses numeric validation
& Verbose flag and consolidated validation chain
& Yes & Yes\\
\midrule
\multicolumn{5}{@{}l}{\textbf{Joint successes, first proposal}} & \textbf{3/5} & \textbf{4/5}\\
\multicolumn{5}{@{}l}{\textbf{Joint successes, final accepted feature implementation}} & \textbf{3/5} & \textbf{5/5}\\
\multicolumn{5}{@{}l}{Additional proposals / secure-code fallbacks} & 0 / 0 & 1 / 0\\
\bottomrule
\end{tabularx}
\par\smallskip
\begin{minipage}{\textwidth}\scriptsize
$^{a}$ ZIP output is missing. The saved boundary check passes.
$^{b}$ User-specified remedy, locally implemented after two unsuccessful model repairs.
$^{c}$ An unimported \code{pickle} name leaves the feature and boundary outcomes unknown. The gate's retry resolves this in Memory+Gate.
\end{minipage}

\end{table*}

\subsection{Repair Contracts and Independent Scope}
With input from our Nostrum Biodiscovery collaborator, we selected RFdiffusion~\cite{rfdiffusion}, ProteinMPNN~\cite{proteinmpnn}, haddock3~\cite{haddock3}, DockQ~\cite{dockq}, AutoDock-Vina~\cite{vina}, and foldseek~\cite{foldseek}. Semgrep 1.101.0 reports 420 warnings. Source and history review identify five Python repair targets outside those warning locations. Foldseek contributes eleven warnings but no suitable project-owned Python case.

RFdiffusion restricts checkpoint loading with \code{weights_only=True} and \code{safe_globals([slice])} under PyTorch 2.5.1~\cite{pytorchserialization,pythonpickle}. Haddock3 checks resolved extraction destinations; ProteinMPNN restricts record names used in output paths. DockQ closes handles on parser failure, and AutoDock-Vina validates numeric arguments across overwrite branches. These last two are reliability contracts. Four repairs were model-written; RFdiffusion used a developer-specified remedy after two unsuccessful model attempts.

Each original boundary assertion fails before repair and passes afterward with an ordinary-input control. External JSON notes record invariants, rationale, source identity, and checks. We strip comments and docstrings with \code{tokenize} and \code{ast}, preserve the executable AST excluding documentation nodes, and reverify the cleaned starts~\cite{pythontokenize,pythonast}.

\subsection{Equal-Budget Feature Evolution}
The tasks extend behavior near these contracts (Table~\ref{tab:bio}), including legacy checkpoint compatibility, archive aliases, hierarchical names, in-memory pose output, and lazy parsing. This purposive stress sample has ten tasks but five repository clusters. No SusVibes continuation enters the evaluation.

We retain each default-mode check and add a check for the new mode: an approved hierarchy or Namespace must not fail merely because the old contract excluded it. Before generation, each repaired start passes its old checks; an evaluator-written reference passes all checks; and a weakened reference retains the feature but triggers a boundary assertion. Reference code is withheld from the model.

An oracle audit removed undocumented requirements to retain Namespace representation and avoid consuming model numbers before parsing. All four conditions of the three affected tasks were rerun. An adapter audit allowed the supplied \code{FilePath = Path} alias without changing generated code and replayed all four archive-alias conditions. The comparison retains 24 unchanged cells, twelve reruns, and four replays. Original runs are archived; prompts, security assertions, and retry budgets were unchanged.

DeepSeek-V3.2 at temperature zero receives the cleaned source and feature request in all four conditions:
\begin{itemize}
\item \textbf{B0: Functional gate.} Application, ordinary, and feature failures can trigger one retry.
\item \textbf{B1: Memory.} B0 plus the repair note and a prospective clarification of the extended contract.
\item \textbf{B2: Boundary gate.} B0 plus saved and extended boundary checks in retry feedback.
\item \textbf{B3: Both.} Memory and boundary feedback.
\end{itemize}
All cells have the same token budget and maximum of two proposals. B0/B1 boundary checks execute only after generation ends; their outcomes cannot affect retries. Checks use pinned offline containers with read-only source and evaluator mounts. All 40 cells are retained, without outcome-dependent stopping. The earlier unequal-budget comparison remains separate archival evidence.

FuncPass requires ordinary and feature checks; SecPass requires saved and extended boundary checks; JointPass requires both. We report first and final proposals separately. Execution errors are unknown boundary outcomes, not violations. A rejected final proposal leaves the repaired start in place; this fallback is not feature delivery.

Table~\ref{tab:bio} separates first and final outcomes. B0 accepts three violating proposals: one unrestricted checkpoint load and two portable backslash-name violations. The latter are contract failures, not demonstrated Linux path escapes. B1 accepts one proposal that relaxes the default simple-name restriction. B2 and B3 accept no observed violation, with seven and one fallbacks, respectively. B0 has one unknown outcome.

The five additional joint successes of B3 over B0 comprise one restricted-loading success, one portable-naming success, and three feature or execution fixes. Figure~\ref{fig:bio-memory} shows the deserialization contrast. The combined condition has the highest joint count in this cohort; five selected repository clusters do not establish a population effect.

\begin{figure}[t]
\centering
% Code excerpt from silent_fix_feature_v1/cases/haddock_extract_symlink/starting_code.py.
% The repair-note fields below are paraphrased for display, not literal JSON.
\fbox{\begin{minipage}{0.94\columnwidth}
\footnotesize
\textbf{1. Silent Security Repair in Code}\par
Comments stripped.\par\smallskip
{\ttfamily
 target\_path = dest / member.name\par
 target\_resolved = target\_path.resolve()\par
 if dest not in target\_resolved.parents \textbackslash\par
 \hspace*{1em}and dest != target\_resolved:\par
 \hspace*{2em}continue\par}
\smallskip\hrule\smallskip
\textbf{2. TRACER External Repair Note}\par
Summary of the saved JSON note.\par
\textbf{Target symbol.} \code{extract_files_flat}\par
\textbf{Security rule.} Final extraction paths must stay inside the destination directory after resolving symlinks.\par
\textbf{Linked boundary check.} Reject a symlink escape and accept safe archive members.\par
\smallskip\hrule\smallskip
\textbf{3. Subsequent Feature Evolution}\par
\textbf{Request.} Add ZIP/TAR and nested-path support.\par
\textbf{No Memory.} Fails the feature check.\par
\textbf{Memory + Gate.} Passes ordinary behavior, the new feature check, and the saved security boundary check.
\end{minipage}}

\caption{An observed RFdiffusion regression under equal retry budgets. Both final proposals pass ordinary and legacy-feature checks; B0 enables unrestricted deserialization, while B3 retains restricted loading and passes both boundary checks.}
\label{fig:bio-memory}
\end{figure}

\rqfinding{RQ5}{Can repair memory and boundary gates preserve repaired properties while subsequent feature changes succeed?}{With the same one-retry limit, joint success is 4/10 with functional feedback alone, 8/10 with memory, 3/10 with a boundary gate, and 9/10 with both. B3 rejects its remaining contract-violating proposal; fallback is not feature success.}

\section{Lessons for SE Teams and Threats to Validity}
\label{sec:threats}
\subsection{Practical Lessons}
\noindent\textbf{Context dilution.} C2 outperforms C3 with both investigators. A relevant explanation may connect an assumption to a current-code obligation more directly than combined structural and textual history. Context dilution is a plausible explanation, not an isolated causal finding, since context length and attention are uncontrolled. Teams should retrieve history to resolve a specific hypothesis rather than maximize context volume.

\noindent\textbf{Verifiable boundary obligations.} SAST supplies reproducible entry points, while patch-selected symbols expose changes without warnings. A missing sibling guard or recorded trust assumption becomes actionable when translated into a check of current inputs, consumers, or object lifecycle. Preserve the link from the historical observation through the hypothesis to its validating evidence.

\noindent\textbf{Validation and maintenance gates.} Reserve final validation capacity, retain rejected patches and reasons, and reuse results only for identical candidates. Recognized terminal summaries distinguish completed checks from interrupted tests. Persistent memory carries the accepted invariant and independently owned check into maintenance, where the gate returns concrete feedback or preserves the secure start. Report feature delivery and fallback preservation separately.

\subsection{Threats to Validity and Scope}
\emph{Construct and internal validity.} Reference-patch overlap measures review opportunity, with broad classes and incidental edits potentially overstating semantic localization. Task tests cover specified contracts, not whole-repository security or unrelated alerts. Test adequacy and overfitting remain limits~\cite{repairsurvey,overfitting}, and recorded rationale need not faithfully describe internal reasoning~\cite{turpin,lanham}. Adaptive development changes several mechanisms together. C0 receives patch-selected targets, C1--C3 add audit cues without a cue-only control, and temporal evidence excludes upstream ancestry. Resumed execution retains earlier revisions, with paired original/candidate evaluations after environment recovery. The previously protected stratum is separately reported but is not an untouched prospective replication.

\emph{External and domain-study validity.} Selection favors Python outputs with publisher-reported functional success. Missing complete outcomes for 42 admissions may be nonrandom, and shared tasks create dependence addressed by task-cluster intervals. Few conversions limit precision. The fifteen semantic reviews are purposive and single-reviewer. Five component repairs in six structural-biology repositories are researcher-run with industry input, not independently confirmed CVEs or a practitioner deployment. Memory+Gate includes a retry absent from No Memory, so its gain cannot be attributed to memory alone. Source/check hashes establish traceability rather than semantic correctness. Broader repository coverage and independent practitioner validation remain future work.

\section{Related Work}
SWE-bench measures issue resolution~\cite{swebench}, whereas SusVibes separates functional and security outcomes~\cite{susvibes,susvibesartifact}. CVEfixes links vulnerabilities to fixing histories, and Vul4J provides reproducible vulnerabilities and triggering tests~\cite{cvefixes,vul4j}. SWE-agent studies repository interaction interfaces~\cite{sweagent}, while ReAct supplies interleaved explanation/action traces~\cite{react}. For security investigation, PAGENT combines static guidance with dynamic profiling~\cite{pagent}, IRIS infers taint specifications and filters candidate paths~\cite{iris}, LLMDFA couples model reasoning with parsing and feasibility checks~\cite{llmdfa}, and RepoAudit explores data flows using memory and validation~\cite{repoaudit}. Our comparison asks what development evidence adds to a common current-source repair protocol.

Historical analysis operates on complementary units. ViolationTracker reconstructs static-warning histories~\cite{violationtracker}, CoLeFunDa represents function changes to identify silent fixes~\cite{colefunda}, KNighter synthesizes checkers from historical patches~\cite{knighter}, and FixMiner extracts reusable repair patterns~\cite{fixminer}. GumTree provides structural edit matching~\cite{gumtree,gumtree3}. Analyzer-guided repair includes InferFix~\cite{inferfix}, while ChatRepair uses conversational test feedback~\cite{chatrepair} and RepairAgent autonomously inspects, edits, and tests code~\cite{repairagent}. These mechanisms motivate tracking losses between available evidence, supported diagnosis, and retained repair.

For subsequent maintenance, Reflexion retains verbal feedback in episodic memory~\cite{reflexion}, Self-Refine iterates through model-generated feedback~\cite{selfrefine}, and POSTURA describes a persistent security graph~\cite{postura}. The exploratory \emph{Undoing the Fix} report studies deliberately elicited regressions under an AI-control setting~\cite{undoingfix}. Our benign feature-evolution study instead connects an earlier repair invariant to independently owned executable checks, distinguishing first proposals, feedback-assisted acceptance, and secure fallback.

\section{Conclusion}
\tracer{} links change-informed planning, temporal AST evolution, recorded rationale, and guarded repair to persistent maintenance evidence. Across 3,096 cells, C2 achieves the highest shared-cohort SecPass counts with both investigators, reaching 76/345 with DeepSeek and 75/345 with Gemini from 68 original passes. Stage-level losses explain why broad review coverage produces modest repair gains. In structural-biology components, memory with a one-retry gate completes five feature evolutions while preserving saved boundaries, compared with three under unguarded evolution. The portable skill and archived evidence support auditable security review and maintenance.

\section*{Data Availability}
The skill, analysis helpers, derived per-cell outcomes, and tables are archived on Zenodo~\cite{tracerartifact} and maintained at \url{https://github.com/pkrajput/TRACER_skill}. SusVibes provides the benchmark inputs~\cite{susvibesartifact}. The archive supports outcome recomputation, but excludes raw investigator transcripts, complete evaluator logs, and container exports.

\section*{Acknowledgements}
We thank Google for Gemini API credits and Nostrum Biodiscovery for verification assistance and feedback on the work's practical relevance.

